\]
and $\pi=16\arctan(1/5)-4\arctan(1/239)$.
Equation~\eqref{eq:logcertificate} follows by integrating the geometric
series for $2/(1-z^2)$ and bounding its positive tail.
Equation~\eqref{eq:atancertificate} is the alternating-series estimate;
the reductions follow from the tangent addition formula with their
stated ranges fixing the branch. The Machin identity follows by the
same addition formula and the bounds placing its value in $(3,4)$.

The source uses $m=64$ for logarithms, $m=80$ for arctangents, and
square-root brackets of width $2^{-144}$. The independent checker
retains the same series lengths and uses outward-rounded dyadic
arithmetic at \emph{192 bits}, including 192-bit square-root brackets.
This changes the implementation precision, not any mathematical constant.

Here is the finite checker contract, intended to be proved once in
Lean. An interval $(l,u)$ represents $[l/2^{192},u/2^{192}]$, with
$l,u\in\Z$ and $l\le u$. Addition and subtraction use endpoint arithmetic;
multiplication encloses the four endpoint products, rounded outward.
Reciprocal reverses the endpoints and is permitted only when zero is
excluded. A nonnegative square root is enclosed by integers whose
squares bracket the appropriately scaled rational endpoint. Each series
is evaluated using these interval operations and enlarged by the
explicit remainder bound above. Monotonicity allows logarithms and
arctangents to be evaluated separately at the two rational endpoints.
Induction over these operations proves enclosure soundness. All branch
choices must be justified by exact endpoint comparisons.

The supplied \path{certificates/interval.py} implements these rules
using integer arithmetic and rational numbers. Decimal output is used
only for display. The complete accepted interval endpoints are stored
in \path{certificates/verification_results.json}. Executing Python is
not a kernel proof; a Lean checker must prove this soundness contract
and replay the finite data or an equivalent exact computation.
\gid{Zeta5.Certificate.interval_soundness}

\subsection{The half-line potential inequality}
For $t>0$, direct integration in \eqref{eq:Vdef} gives
\begin{equation}\label{eq:Vexplicit}
\begin{split}
 V(t)={}&\log(1+t)-6\alpha\log(t+\alpha^2)-2+12\alpha\\
 &+2\sqrt t\left(\pi+\arctan\frac1{\sqrt t}
                         -6\arctan\frac\alpha{\sqrt t}\right).
\end{split}
\end{equation}
At zero, $V(0)=-12\alpha\log\alpha-2+12\alpha$. If
$\Phi(y)=V(y^2)$, then
\[
 \Phi'(y)=2(\pi+\arctan(1/y)-6\arctan(\alpha/y)),\qquad
 \Phi''(y)=-\frac2{1+y^2}+\frac{12\alpha}{\alpha^2+y^2}.
\]
The second derivative changes sign once, at $y^2=711/880$.
Moreover $\Phi'(0+)=-3\pi$ and $\Phi'(+\infty)=2\pi$.
It follows that $V$ has exactly one minimum: $\Phi'$ first increases
from a negative value and then decreases to a strictly positive limit.
Let
\[
 q_- =\frac{59205077}{10^{10}},\qquad
 q_+ =\frac{59205079}{10^{10}}.
\]
The interval evaluation gives
$\Phi'(\sqrt{q_-})<0<\Phi'(\sqrt{q_+})$, so the minimizing $t$ belongs
to $[q_-,q_+]$. On that interval a lower bound is
\begin{equation}\label{eq:Vstar}
\begin{split}
 V_*={}&\log(1+q_-)-6\alpha\log(q_++\alpha^2)-2+12\alpha\\
 &+2\sqrt{q_+}\left(\pi+\arctan\frac1{\sqrt{q_+}}
                         -6\arctan\frac\alpha{\sqrt{q_-}}\right).
\end{split}
\end{equation}
The parenthesis is negative, as the same interval computation verifies.
Bounding each logarithm and arctangent in its monotone direction and
using this negative lower bound justifies the factor $\sqrt{q_+}$.
In particular $V_*$ is a lower bound for the global minimum of $V$.

The nested-support potential $U^\rho$ is nonincreasing up to $a_1$,
constant on $[a_1,b_1]$, and nondecreasing thereafter. Thus its maximum
on any interval is at an endpoint. For $[l,r]\subseteq[0,2]$ define
\begin{equation}\label{eq:Bcell}
 B(l,r)=2\max(U^\rho(l),U^\rho(r))-
 \begin{cases}
 V(r),&r\le q_-,\\
 V(l),&l\ge q_+,\\
 V_*,&\text{otherwise}.
 \end{cases}
\end{equation}
The preceding monotonicity proves
$2U^\rho(t)-V(t)\le B(l,r)$ for $l\le t\le r$.

Define the ordered rational endpoints
\[
 A_0=0,\quad A_j=a_{17-j}\ (1\le j\le16),\quad
 A_{17}=q_-,\ A_{18}=q_+,\quad
 A_{18+j}=b_j\ (1\le j\le16),\quad A_{35}=2.
\]
A row $(j,d,k)$ in Table~\ref{tab:partition} denotes the closed cell
\begin{equation}\label{eq:cell}
 \left[A_j+(A_{j+1}-A_j)\frac{k}{2^d},
       A_j+(A_{j+1}-A_j)\frac{k+1}{2^d}\right].
\end{equation}
Every integer in a displayed inclusive range is used. The exact endpoint
checks show that the 684 resulting cells cover $[0,2]$ in order, with
adjacent endpoints equal and with no overlap of interiors.
\begingroup\small
\begin{longtable}{ccl}
\caption{The rational partition (source Table 2).}\label{tab:partition}\\
\toprule $j$ & $d$ & $k$ (inclusive ranges)\\\midrule\endfirsthead
\multicolumn{3}{c}{Table \thetable\ continued}\\
\toprule $j$ & $d$ & $k$ (inclusive ranges)\\\midrule\endhead
\midrule\multicolumn{3}{r}{Continued on the next page}\\\endfoot
\bottomrule\endlastfoot
0 & 1 & $0$\\
0 & 2 & $2,3$\\
1 & 0 & $0$\\
2 & 0 & $0$\\
3 & 0 & $0$\\
4 & 0 & $0$\\
5 & 0 & $0$\\
6 & 0 & $0$\\
7 & 0 & $0$\\
8 & 0 & $0$\\
9 & 1 & $0,1$\\
10 & 1 & $0,1$\\
11 & 1 & $0,1$\\
12 & 1 & $0$\\
12 & 2 & $2,3$\\
13 & 2 & $0,\ldots,3$\\
14 & 2 & $0,\ldots,3$\\
15 & 1 & $0$\\
15 & 2 & $2,3$\\
16 & 0 & $0$\\
17 & 0 & $0$\\
18 & 0 & $0$\\
19 & 1 & $1$\\
19 & 3 & $0,\quad 2,3$\\
19 & 4 & $2,3$\\
20 & 2 & $3$\\
20 & 3 & $4,5$\\
20 & 4 & $0,1,\quad 6,7$\\
20 & 5 & $4,\ldots,11$\\
21 & 2 & $3$\\
21 & 3 & $5$\\
21 & 4 & $0,\quad 7,\ldots,9$\\
21 & 5 & $2,\ldots,4,\quad 10,\ldots,13$\\
21 & 6 & $10,\ldots,19$\\
22 & 3 & $5,\ldots,7$\\
22 & 4 & $0,\quad 8,9$\\
22 & 5 & $2,3,\quad 12,\ldots,15$\\
22 & 6 & $8,\ldots,23$\\
23 & 3 & $6,7$\\
23 & 4 & $9,\ldots,11$\\
23 & 5 & $0,\ldots,2,\quad 14,\ldots,17$\\
23 & 6 & $6,\ldots,13,\quad 19,\ldots,27$\\
23 & 7 & $28,\ldots,37$\\
24 & 3 & $6,7$\\
24 & 4 & $10,11$\\
24 & 5 & $0,\ldots,2,\quad 15,\ldots,19$\\
24 & 6 & $6,\ldots,11,\quad 22,\ldots,29$\\
24 & 7 & $24,\ldots,43$\\
25 & 3 & $7$\\
25 & 4 & $10,\ldots,13$\\
25 & 5 & $0,\ldots,2,\quad 16,\ldots,19$\\
25 & 6 & $6,\ldots,11,\quad 24,\ldots,31$\\
25 & 7 & $24,\ldots,47$\\
26 & 3 & $7$\\
26 & 4 & $11,\ldots,13$\\
26 & 5 & $0,\ldots,2,\quad 17,\ldots,21$\\
26 & 6 & $6,\ldots,11,\quad 25,\ldots,33$\\
26 & 7 & $24,\ldots,49$\\
27 & 4 & $11,\ldots,15$\\
27 & 5 & $0,\ldots,2,\quad 17,\ldots,21$\\
27 & 6 & $6,\ldots,11,\quad 26,\ldots,33$\\
27 & 7 & $24,\ldots,51$\\
28 & 4 & $12,\ldots,15$\\
28 & 5 & $0,\ldots,2,\quad 18,\ldots,23$\\
28 & 6 & $6,\ldots,11,\quad 27,\ldots,35$\\
28 & 7 & $24,\ldots,53$\\
29 & 4 & $13,\ldots,15$\\
29 & 5 & $0,\ldots,2,\quad 19,\ldots,25$\\
29 & 6 & $6,\ldots,12,\quad 27,\ldots,37$\\
29 & 7 & $26,\ldots,53$\\
30 & 4 & $13,\ldots,15$\\
30 & 5 & $0,\ldots,2,\quad 20,\ldots,25$\\
30 & 6 & $6,\ldots,14,\quad 26,\ldots,39$\\
30 & 7 & $30,\ldots,51$\\
31 & 4 & $14,15$\\
31 & 5 & $0,\ldots,2,\quad 20,\ldots,27$\\
31 & 6 & $6,\ldots,39$\\
32 & 4 & $15$\\
32 & 5 & $0,\ldots,4,\quad 19,\ldots,29$\\
32 & 6 & $10,\ldots,37$\\
33 & 4 & $0,\quad 15$\\
33 & 5 & $2,\ldots,29$\\
34 & 2 & $2,3$\\
34 & 3 & $3$\\
34 & 4 & $3,\ldots,5$\\
34 & 5 & $4,5$\\
34 & 6 & $5,\ldots,7$\\
34 & 7 & $6,\ldots,9$\\
34 & 8 & $7,\ldots,11$\\
34 & 9 & $5,\ldots,13$\\
34 & 10 & $0,\ldots,9$\\
\end{longtable}
\endgroup
For every listed cell the outward-rounded computation proves
\begin{equation}\label{eq:cellcertificate}
 B(l,r)<-\frac{6645002}{10^6}<-\frac{1329}{200}.
\end{equation}
For reproducibility, the largest certified upper endpoint occurred in
cell $(j,d,k)=(30,6,26)$ and is smaller than
$-664500268/10^8=-6.64500268$. This is an upper bound on a cell bound,
not a computation of the actual supremum of $2U^\rho-V$.
For $t\ge2$, \eqref{eq:fieldtail} is decreasing and is already less than
$M_0$ at $2$. This proves \eqref{eq:potentialbound} on the entire
half-line.
\gid{Zeta5.Certificate.potential_cells}

\subsection{Energy, scalar constant, and the optional norm constant}
The same interval rules applied to \eqref{eq:energyexplicit} and
\eqref{eq:Cstar} give
\begin{align}
 -\frac{2126593445148}{10^{12}}&<I(\rho)
                         <-\frac{2126593445147}{10^{12}},\label{eq:energycert}\\
 \frac{2653035990340}{10^{12}}&<C_*
                         <\frac{2653035990341}{10^{12}}.\label{eq:Cstarcert}
\end{align}
Since $\lambda M_0=-49173/8000$, it follows that
\[
 \lambda M_0-I(\rho)+C_*<
 -\frac{1366995564511}{10^{12}}<-\frac{2733991}{2000000}=U.
\]
This completes Proposition~\ref{prop:comparison}.
For the constant $\kappa$ in \eqref{eq:kappaexplicit} below, the checker
also proves
\begin{equation}\label{eq:kappacert}
 \frac{2773171335773}{200000000000}<\kappa
 <\frac{6932928339433}{500000000000}<\frac{6933}{500}.
\end{equation}
\gid{Zeta5.Certificate.energy_and_norm_constants}

\section{Exact arithmetic integrals and margins}
\label{app:arithmetic}
\subsection{A branch-complete calculation of the inner integral}
For $0\le u<1$, define
\[
 d_0(u)=\min(u,1-u),\qquad
 e(u)=\begin{cases}1,&u\le1/2,\\-1,&u>1/2.\end{cases}
\]
For real $x$ and almost every $z\in(0,1/2)$,
\[
 \ell(x,z)-2x=e(\{x\})
                  \bigl(\boldsymbol1_{z<d_0(\{x\})}-2d_0(\{x\})\bigr).
\]
The exceptional equality case in $z$ is a singleton. Thus, with $f,g,A,B$
as in Section~\ref{sec:normalization},
\begin{align}
 \int_0^{1/2}B^2&=d_0(g)(1-2d_0(g)),\label{eq:BB}\\
 \int_0^{1/2}AB&=e(f)e(g)
  \bigl(\min(d_0(f),d_0(g))-2d_0(f)d_0(g)\bigr).\label{eq:AB}
\end{align}
These follow by integrating indicator functions and their product;
for example the intersection has length $\min(d_0(f),d_0(g))$.

Take the ordered set consisting of $3,20$, and every $x\in(3,20)$ for
which $cx\in\Z$ for some
\begin{equation}\label{eq:breakpoints}
 c\in\{2,2\alpha,2\lambda,2H,4\alpha,2(1-\alpha),2(1+\alpha)\}.
\end{equation}
There are exactly 143 open intervals $(t_i,t_{i+1})$. On each, all
floors, signs, and branches of minima and positive parts are fixed.
The extra possible crossings $d_0(f)=d_0(g)$ and $\tau=\sigma$ are
included in \eqref{eq:breakpoints}. The quadratic coefficient in $R$
cancels because
\[
 2\lambda-12\lambda\alpha+2(H^2-H-\lambda^2)-18\alpha^2+6\alpha=0.
\]
Therefore $R(x)=a_ix+b_i$ on each such open interval, and
\begin{equation}\label{eq:innerpieceintegral}
 \int_3^{20}R(x)x^{-3}\,dx
 =\sum_{i=0}^{142}\left[
 a_i\left(\frac1{t_i}-\frac1{t_{i+1}}\right)
 +\frac{b_i}{2}\left(\frac1{t_i^2}-\frac1{t_{i+1}^2}\right)\right].
\end{equation}
Values at finitely many endpoints do not affect the integrals.

The independent verifier does not merely sample two points and assume
linearity. It forms exact rational polynomials on each branch, verifies
every branch inequality at both endpoints, and computes $\Gamma$ from
its original integral on the strips in $z$ bounded by
$0,d_0(f),d_0(g),1/2$. It separately forms \eqref{eq:RFQ}--\eqref{eq:Qtail}
using \eqref{eq:BB}--\eqref{eq:AB}, checks equality of all polynomial
coefficients, and checks cancellation of degrees greater than one.
Then it integrates the resulting affine polynomial exactly by
\eqref{eq:innerpieceintegral}. The 143 coefficient records are included
in the JSON output. Grouping them by integer unit intervals gives
Table~\ref{tab:innerintegrals}; its sum is \eqref{eq:innerexact}.
\gid{Zeta5.Certificate.inner_integral}
\begin{table}[htbp]\centering\small
\caption{Inner integral on unit intervals (source Table 3).}\label{tab:innerintegrals}
\renewcommand{\arraystretch}{1.65}
